\begin{figure}[ht]
\centering
\begin{tikzpicture}[x=1cm, y=1cm, font=\small]

% Log-log axes
\draw[->, thick] (-0.2, 0) -- (8.5, 0) node[right, font=\footnotesize] {$\log_{10}(\tau / \text{s})$};
\draw[->, thick] (0, -0.2) -- (0, 6.0) node[above, font=\footnotesize] {$\log_{10}(\sigma_y(\tau))$};

% Tick marks on tau axis
\foreach \x in {0,2,4,6,8} {
  \draw[thick] (\x, -0.1) -- (\x, 0.1);
  \node[anchor=north, font=\footnotesize] at (\x, -0.1) {\x};
}
% Tick marks on sigma axis (values are log10)
\foreach \y in {0,1,2,3,4,5} {
  \pgfmathsetmacro{\v}{-12 + \y}
  \draw[thick] (-0.1, \y) -- (0.1, \y);
  \node[anchor=east, font=\footnotesize] at (-0.1, \y) {\v}
  ;
}

% Quartz OCXO curve: white-frequency at short tau, flicker-floor middle, drift at long tau
% Approximate piecewise log-log line
\draw[engblue, very thick]
  (0.5, 5.5) -- (2.0, 3.5)  % white FM slope -1/2
  (2.0, 3.5) -- (4.5, 3.5)  % flicker floor
  (4.5, 3.5) -- (7.5, 5.0); % drift up
\node[engblue, anchor=west, font=\footnotesize] at (7.6, 5.0) {Quartz OCXO};

% Caesium beam: lower white-FM, flicker floor lower, gentle drift
\draw[engred, very thick]
  (0.5, 4.5) -- (3.0, 2.0)
  (3.0, 2.0) -- (5.0, 2.0)
  (5.0, 2.0) -- (7.5, 3.0);
\node[engred, anchor=west, font=\footnotesize] at (7.6, 3.0) {Cs beam};

% H-maser: best short and medium term, drift kicks in late
\draw[engblack, very thick]
  (0.5, 3.0) -- (4.0, 0.5)
  (4.0, 0.5) -- (6.0, 0.5)
  (6.0, 0.5) -- (7.5, 1.5);
\node[engblack, anchor=west, font=\footnotesize] at (7.6, 1.5) {H-maser};

% Caesium fountain: lowest floor
\draw[engblue!50!black, very thick, dashed]
  (1.5, 2.0) -- (5.0, -0.5)
  (5.0, -0.5) -- (7.0, -0.5);
\node[engblue!50!black, anchor=west, font=\footnotesize] at (7.1, -0.5) {Cs fountain};

% Annotations for the slope regions
\node[font=\footnotesize, anchor=south] at (1.0, 6.0) {white FM};
\node[font=\footnotesize, anchor=south] at (4.0, 6.0) {flicker};
\node[font=\footnotesize, anchor=south] at (7.0, 6.0) {drift};
\draw[gray, dashed] (2.5, -0.2) -- (2.5, 5.7);
\draw[gray, dashed] (5.5, -0.2) -- (5.5, 5.7);

\end{tikzpicture}
\caption[Allan deviation as a function of averaging time]{Sketch of
the Allan deviation $\sigma_y(\tau)$ versus averaging time $\tau$
for four oscillator classes (log-log axes, illustrative). The
short-$\tau$ regime is dominated by white frequency noise (slope
$-1/2$ on the log-log plot); a flicker-noise floor sets the
medium-$\tau$ minimum; long-$\tau$ behaviour rises again under
slow drift. The Allan deviation distinguishes oscillator types
where a single-number fractional accuracy cannot: two oscillators
with the same one-second stability may differ by orders of
magnitude at one-day averaging.}
\label{fig:vol01:ch06:allan-deviation}
\end{figure}
